\SourceLeaf{59}
\begin{center}
{\large THEOREMATIS FUNDAMENTALIS\par}
\smallskip IN\par
\smallskip{\large DOCTRINA DE RESIDUIS QUADRATICIS\par}
\medskip DEMONSTRATIONES ET AMPLIATIONES NOVAE.\par
\end{center}
\BookRule{20mm}

Theorema fundamentale de residuis quadraticis, quod inter pulcherrimas arithmeticae sublimioris veritates refertur, facile quidem per inductionem detectum, longe vero difficilius demonstratum est. Saepius in hoc genere accidere solet, ut veritatum simplicissimarum, quae scrutatori per inductionem sponte quasi se offerunt, demonstrationes profundissime lateant et post multa demum tentamina irrita, longe forte alia quam qua quaesitae erant via, tandem in lucem protrahi possint. Dein haud raro fit, quam primum una inventa est via, ut \textit{plu}\textit{res} subinde patefiant ad eandem metam perducentes, aliae brevius et magis directe, aliae quasi ex obliquo et a principiis longe diversis exorsae, inter quae et quaestionem propositam vix ullum vinculum suspicatus fuisses. Mirus huiusmodi nexus inter veritates abstrusiores non solum peculiarem quandam venustatem hisce contemplationibus conciliat, sed ideo quoque sedulo investigari atque enodari meretur, quod haud raro nova ipsius scientiae subsidia vel incrementa inde demanant.

Etsi igitur theorema arithmeticum, de quo hic agetur, per curas anteriores, quae quatuor demonstrationes inter se prorsus diversas\begingroup\renewcommand{\thefootnote}{*)}\footnote{Duae expositae sunt in \textit{Disquisitionum Arithmeticarum} Sect. quarta et quinta; tertia in commentatione peculiari (\textit{Commentt. Soc. Gotting. Vol. XVI}), quarta inserta est commentationi: \textit{Summatio quarundam} \textit{serierum singularium} (\textit{Commentt. Recentiores, Vol. I}).}\endgroup suppeditaverunt, plene
\SourceLeaf{60}

absolutum videri possit, tamen denuo ad idem argumentum revertor, duasque alias demonstrationes adiungo, quae novam certe lucem huic rei affundent. Prior quidem tertiae quodammodo affinis est, quod ab eodem lemmate proficiscitur; postea vero iter diversum prosequitur, ita ut merito pro demonstratione nova haberi possit, quae concinnitate ipsa illa tertia si non superior saltem haud inferior videbitur. Contra demonstratio sexta principio plane diverso subtiliori innixa est novumque sistit exemplum mirandi nexus inter veritates arithmeticas primo aspectu longissime ab invicem remotas. Duabus hisce demonstrationibus adiungitur algorithmus novus persimplex ad diiudicandum, utrum numerus integer datus, numeri primi dati residuum quadraticum sit an non-residuum.

Alia adhuc affuit ratio, quae ut novas demonstrationes, novem iam abhinc annos promissas, nunc potissimum promulgarem, effecit. Scilicet quum inde ab anno 1805 theoriam residuorum cubicorum atque biquadraticorum, argumentum longe difficilius, perscrutari coepissem, similem fere fortunam, ac olim in theoria residuorum quadraticorum, expertus sum. Protinus quidem theoremata ea, quae has quaestiones prorsus exhauriunt, et in quibus mira analogia cum theorematibus ad residua quadratica pertinentibus eminet, per inductionem detecta fuerunt, quam primum via idonea quaesita essent: omnes vero conatus, ipsorum demonstrationibus ex omni parte perfectis potiundi, per longum tempus irriti manserunt. Hoc ipsum incitamentum erat, ut demonstrationibus iam cognitis circa residua quadratica alias aliasque addere tantopere studerem, spe fultus, ut ex multis methodis diversis una vel altera ad illustrandum argumentum affine aliquid conferre posset. Quae spes neutiquam vana fuit, laboremque indefessum tandem successus prosperi sequuti sunt. Mox vigiliarum fructus in publicam lucem edere licebit: sed antequam arduum hoc opus aggrediar, semel adhuc ad theoriam residuorum quadraticorum reverti, omnia quae de eadem adhuc supersunt agenda absolvere, atque sic huic arithmeticae sublimioris parti quasi valedicere constitui.
